\section{Local projectors and exact spectral repair}
\label{sec:exact-repair}

Boundary completion leaves a simple regular graph whose nonconstant spectrum
lies within $e/20$ of the desired interval.  We remove this error by switching
edges.  Two estimates make the switches effective: the spectral mass near each
edge is small at every vertex, and the bulk resolvent has a prescribed
$2$-by-$2$ block on most edges.  We first derive these estimates, then prove a
repair theorem that uses only these two properties and spectral confinement.

Throughout this section, put
\[
 q=d-1,\qquad s=2\sqrt q,\qquad a=10^{-4},\qquad
 e=n^{-2/3+a}.
\]
All operators initially defined on $\mathbf 1^\perp$ are extended by zero on
$\operatorname{span}\{\mathbf 1\}$.  Choose an orthonormal eigenbasis
$v_2,\ldots,v_n$ of $\mathbf 1^\perp$ for the adjacency matrix $A_G$, and write
\[
 A_Gv_i=\lambda_i v_i,\qquad x_i^\sigma=s-\sigma\lambda_i
 \quad (\sigma\in\{+1,-1\}).
\]
The graph supplied by the pairing construction has $x_i^\sigma>-e/20$.

\subsection{Extraction of the local spectral estimates}

We extract the two local estimates from different features of the reference.
For the spectral mass bound, a bounded diagonal derivative in $h$ gives a
resolvent difference of order $\sqrt w$, since $z(h)-s$ is quadratic in
$h_c-h$, where $h_c=q^{-1/2}$.  A difference quotient then bounds the
squared resolvent by $O(w^{-1/2})$, and hence the mass at slack at most $w$
by $O(w^{3/2})$.  The typical edge blocks will instead use the path
information retained at the separation stop.

\begin{lemma}[Projector estimate from diagonal comparison]
\label{lem:repair-projector}
Assume $x_i^\sigma>-e/20$ for all nonconstant indices $i$ and both signs.
Suppose the final resolvents
$R_\sigma(z)=(z-\sigma A_G)^{-1}|_{\mathbf 1^\perp}$ at the primary points
$z=z(h)$ satisfy the entry comparison
\[
 \max_p\left|
 R_\sigma(z(h))_{pp}-m_h\widehat P_{\mathrm{fin},pp}^\sigma(h)
 \right|\le C B_B,
 \qquad B_B=o(\sqrt e).
\]
Suppose also that the same frozen reference is defined throughout the primary
range and obeys
\[
 \sup_{h,p}\left|
 \frac{d}{dh}\bigl(m_h\widehat P_{\mathrm{fin},pp}^\sigma(h)\bigr)
 \right|\le C.
\]
Assume the primary grid has bounded geometric spacing, as in the construction.
Then, uniformly in $p$ and $\sigma$,
\begin{equation}
 \sum_{i:\,x_i^\sigma\le w}|v_i(p)|^2\le Cw^{3/2}
 \qquad (2e\le w\le c_d).
 \label{eq:original-61}
\end{equation}
The constant can be enlarged to make the same assertion valid for all larger
$w$ in the bounded spectral range.
\end{lemma}

\begin{proof}
List the primary gaps increasingly as $t_0<t_1<\cdots<t_N$, with
$1<r_{\min}\le t_{j+1}/t_j\le r_{\max}$ for fixed constants
$r_{\min},r_{\max}$.
The closest gap is less than $e/20$, and the largest gap is bounded below by
a positive constant.  For $2e\le w\le c_d$, choose the least $j$ with
$t_j\ge w$, and set $w_-=t_j$, $w_p=t_{j+1}$.
Choose $c_d$ small enough that this successor always exists.  Then
$w\le w_-<w_p\le r_{\max}^2w$ and
$w_p-w_-\ge(r_{\min}-1)w$.
In particular both resolvents are positive.  On each eigenspace,
\[
 \frac{1}{(x_i^\sigma+w_p)^2}
 \le
 \frac{(x_i^\sigma+w_-)^{-1}-(x_i^\sigma+w_p)^{-1}}
 {w_p-w_-}.
\]
The identity $z(h)-s=(1-\sqrt q\,h)^2/h$ gives
$|h-h_-|=O(\sqrt w)$ for the corresponding primary parameters.
The diagonal comparison and its derivative bound therefore give
\begin{equation}
 \bigl[R_\sigma(s+w_p)^2\bigr]_{pp}
 \le
 \frac{[R_\sigma(s+w_-)-R_\sigma(s+w_p)]_{pp}}{w_p-w_-}
 \le \frac{C}{\sqrt w}.
 \label{eq:original-62}
\end{equation}
Here the entry error contributes $O(B_B/w)=o(w^{-1/2})$ uniformly because
$w\ge2e$.  If $x_i^\sigma\le w$, the positive denominator
$x_i^\sigma+w_p$ is at most $Cw$.  Multiplying the preceding estimate by
$Cw^2$ proves \eqref{eq:original-61}.  For $w$ bounded below by $c_d$, use
$\sum_i|v_i(p)|^2\le1$ and enlarge the constant.
\end{proof}

\subsubsection*{Verifying the reference derivative}
To apply Lemma~\ref{lem:repair-projector}, we need a derivative bound for the
same reference at every primary parameter.  Let $S_0$ be the boundary at the
separation stop and $l_0=|S_0|$.
Fix the entire matching added after this stop, independently of $h$.
On the contrast space of $S_0$, let $P_f$ be its involution, and put
\[
 H_{f0}=-\sigma hP_f,\qquad
 H_f=H_{f0}(I+\widehat D_0H_{f0})^{-1}.
\]
Lemma~\ref{lem:boundary-smooth-reference} and the exact-reference construction
following it, in particular \eqref{eq:original-40}--\eqref{eq:original-43},
imply uniformly over the primary range
\[
 \|\widehat D_0\|\le C/R_0,\qquad
 \|\widehat D_0\Pi e_t\|+
 \|\partial_h\widehat D_0\Pi e_t\|\le n^{-c},
 \qquad \|\partial_h\widehat D_0\|\le C(1+L)^4
\]
for some fixed $c>0$, after decreasing $c$ to absorb logarithms.  For the last
bound, differentiate the finite stop sum defining $\widehat P_0$ term by
term, keeping its coefficients $Q_j$ and its taper fixed.  By
\eqref{eq:original-7}, the derivative of its $j$th term has norm at most
$Cj(1+j)^2/h$.  The primary interval is bounded away from zero, so
summing over $j<L$ gives $C(1+L)^4$; compression by $V_0$ cannot
increase this norm.  Consequently $H_f$ is bounded and $H_f'$ has at most
polylogarithmic norm.  Moreover,
\[
 (H_f-H_{f0})\Pi e_t
 =-H_{f0}(I+\widehat D_0H_{f0})^{-1}
       \widehat D_0H_{f0}\Pi e_t
\]
has power-small norm, as does its derivative: $H_{f0}\Pi e_t$ is a bounded
multiple of the centered coordinate of the partner of $t$.

For a global vertex $p$, write the stop row as
$p_{p,0}=c_p\Pi e_{t(p)}+\tau_p$, allowing the first term to be absent.
Equation \eqref{eq:original-41} gives power-small norms for $\tau_p$ and
$\tau_p'$.  Also $|c_p|+|c_p'|\le C$: indeed $c_p=(\sigma h)^r$ for a path
length $r$, and $\sup_{r\ge0}rh^{r-1}$ is bounded on the fixed primary range.
The diagonal correction in the frozen reference is $-p_{p,0}^TH_fp_{p,0}$.
Its terms containing a tail have bounded, in fact power-small, derivatives.
For its remaining term, the ideal involution has no fixed point and hence
\[
 (\Pi e_t)^TH_{f0}\Pi e_t=\frac{\sigma h}{l_0}.
\]
The preceding coordinate estimate controls the replacement of $H_{f0}$ by
$H_f$, including its derivative.  Finally the stop diagonal and its derivative
are bounded by the return-path estimates and the virtual-path estimates.
Since $m_h$ and $m_h'$ are bounded, the derivative hypothesis of
Lemma~\ref{lem:repair-projector} follows.  The final entry comparison has error
$O(B_B)$, and the chosen exponents give
\[
 B_B=\frac{n^{.40a}}{\sqrt{l_2}}
     \asymp n^{-1/3+.325a+(g+r)/2}=o(\sqrt e).
\]

\subsubsection*{The bulk resolvent and its typical edge blocks}
Set
\[
 t_*=100a,\qquad b_*=10a,\qquad
 W_b=n^{-2/3+t_*},\qquad M=\lceil n^{1/3+b_*}\rceil.
\]
For a cutoff $w_B^\sigma\in[W_b,2W_b]$ avoiding the signed slacks
$\{x_i^\sigma:2\le i\le n\}$, define
\[
 E_0^\sigma=\operatorname{span}\{v_i:x_i^\sigma<w_B^\sigma\},
 \qquad
 R_B^\sigma=\sum_{x_i^\sigma\ge w_B^\sigma}
            \frac{v_iv_i^T}{x_i^\sigma}.
\]
The reciprocal is positive on its bulk subspace and zero on the complementary
near subspace.  The comparison block is
\[
 h_c=q^{-1/2},\qquad m_c=\frac{h_c}{1-h_c^2},\qquad
 P_{\mathrm{tree}}^\sigma
 =m_c\begin{pmatrix}1&\sigma h_c\\ \sigma h_c&1\end{pmatrix}.
\]

\begin{lemma}[Bulk estimates and typical edges]
\label{lem:repair-bulk}
The projector estimate \eqref{eq:original-61} implies
\begin{equation}
 \sup_p\bigl[(R_B^\sigma)^2\bigr]_{pp}
 \le \frac{C}{\sqrt{W_b}}.
 \label{eq:original-63}
\end{equation}
If $w_p\asymp W_b$ is a primary gap, every entry of
$R_B^\sigma-R_\sigma(s+w_p)$ is $O(\sqrt{W_b})$.

For the typical-edge assertion below, the additional inputs are the following
specific conclusions at the separation stop:
\begin{enumerate}
\item Only $O(l_0)$ edges of $G$ were added after the stop, with
$l_0\asymp n^{2/3+.15a}$.  At most $n^{.08+o(1)}$ stop vertices have a
nonempty actual nonbacktracking return of length at most
$r_n=\lfloor .08\log_q n\rfloor+2$.
\item The number of actual nonbacktracking paths of any length $k<L$ between
fixed endpoints is bounded by a constant.  The stop reference is the weighted
sum of these paths with coefficient $\varphi(k/L)(\sigma h)^k$, up to a
matrix whose entries are $o(n^{-.04})$.  The taper is bounded and equals one
for $k\le r_n$, and $h\le h_c$.
\item The final reference has the form
$\widehat P_{\mathrm{fin}}=\widehat P_0-
 \widehat P_0V_0H_fV_0^T\widehat P_0$, where $V_0$ is an isometry from an
$(l_0-1)$-dimensional space, $\|H_f\|\le C$, and
$\|\widehat P_0\|\le C(1+L)^3$.
At the comparison primary points the final resolvent has entry error
$O(B_B)=o(n^{-.04})$ relative to $m_h\widehat P_{\mathrm{fin}}$.
\end{enumerate}
The return count, multiplicity and virtual-path assertions are consequences of
Lemmas~\ref{lemma:structure} and \ref{lemma:path-majorants}; the norm bound is
Lemma~\ref{lemma:pencil-coefficients}; the reference representation and final
error are furnished by Section~\ref{sec:boundary-completion} and
Proposition~\ref{prop:boundary-completion}.  All inputs are uniform in the
primary points used below.
For each pair of cutoffs as above, all but a fraction $O(n^{-.1})$ of the oriented edges
$(p,p')$ have, for both signs, their endpoint block of $R_B^\sigma$ within
$n^{-.024}$ of $P_{\mathrm{tree}}^\sigma$.  The estimates are uniform in the
cutoffs; their exceptional sets may depend on the cutoffs.
\end{lemma}

\begin{proof}
Partition the bulk slacks into dyadic intervals
$[2^jW_b,2^{j+1}W_b)$.  Their contributions to
$[(R_B^\sigma)^2]_{pp}$ are at most $C(2^jW_b)^{-1/2}$ by
\eqref{eq:original-61}.  Their sum, followed by the bounded contribution from
slacks bounded below by a constant, proves \eqref{eq:original-63}.

On the excluded subspace the diagonal mass of $R_\sigma(s+w_p)$ is at most
$CW_b^{3/2}/W_b=C\sqrt{W_b}$, since $x_i^\sigma>-e/20$ and $W_b\gg e$.
On the included subspace the positive difference of the reciprocal weights is
\[
 \frac1{x_i^\sigma}-\frac1{x_i^\sigma+w_p}
 \le \frac{CW_b}{(x_i^\sigma)^2}.
\]
Its diagonal mass is $O(\sqrt{W_b})$ by \eqref{eq:original-63}.
Cauchy--Schwarz, applied separately to these two positive operators, gives the
entry comparison.

For the typical-edge assertion, the edges added after the separation stop
contribute only $O(l_0)$ exceptions.  Set
$r_n=\lfloor .08\log_q n\rfloor+2=.08\log_q n+O(1)$.
At that stop, the global diagram estimate for one return shows that only
$n^{.08+o(1)}$ vertices have a nonempty nonbacktracking return of length at most
$r_n$.  For an old edge with neither endpoint in this set, the direct edge is
the only nonbacktracking path between its endpoints in the short range
$1\le k\le r_n-1$: an alternative, closed with the direct edge, would give a
nonempty endpoint return of length at most $k+1\le r_n$, after stripping any
terminal reversal.  The diagonal short-path term is the identity.
Bounded fixed-endpoint multiplicity bounds the remaining weighted actual
paths, whose lengths begin at $r_n$, by
\[
 C\sum_{k=r_n}^{L-1}q^{-k/2}
 \le C_dq^{-r_n/2}=n^{-.04+o(1)}.
\]
The virtual and projection terms are smaller.  Thus the stop reference has
the normalized endpoint block
\[
 \begin{pmatrix}1&\sigma h\\ \sigma h&1\end{pmatrix}
 +O(n^{-.04+o(1)}).
\]

The frozen correction is $-p_{p,0}^TH_fp_{p',0}$, with $\|H_f\|\le C$.
The coefficient bound \eqref{eq:original-7} gives
\[
 \sum_p\|p_{p,0}\|^2
 \le l_0\|\widehat P_0\|^2
 \le Cl_0(1+L)^6.
\]
Consequently, outside a fraction at most $n^{-.2}$ of the vertices, the squared
row norms are at most $n^{-.03}$.  Take the union of the exceptional sets for
the two signs at their comparison points.  Fixed degree transfers these
vertex exceptions to the same order of edge exceptions, and the frozen
correction is then $O(n^{-.03})$.  Finally,
$h-h_c=O(\sqrt{W_b})$, the final entry error is $O(B_B)$, and the bulk entry
comparison costs $O(\sqrt{W_b})$.  These errors are smaller than
$n^{-.024}$.  All the exceptional fractions just counted are
$O(n^{-.1})$, as asserted.
\end{proof}

\subsection{A deterministic exact-repair theorem}

\begin{theorem}[Exact repair by simultaneous switches]
\label{thm:exact-repair}
Fix an integer $d\ge3$ and a constant $C_0\ge1$.  Put
\[
 \begin{gathered}
 q=d-1,\qquad s=2\sqrt q,\qquad a=10^{-4},\\
 e=n^{-2/3+a},\qquad W_b=n^{-2/3+100a},\qquad
 M=\lceil n^{1/3+10a}\rceil.
 \end{gathered}
\]
There are constants $n_0,k$, depending only on $d,C_0$, and a deterministic
algorithm with the following property for $n\ge n_0$.

Let $G$ be a simple $d$-regular graph on $n$ vertices.  Let
$v_2,\ldots,v_n$ be an orthonormal eigenbasis of $A_G$ on
$\mathbf1^\perp$, with eigenvalues $\lambda_i$, and put
$x_i^\sigma=s-\sigma\lambda_i$ for $\sigma\in\{+1,-1\}$.
Assume the following three properties.
\begin{enumerate}
\item \emph{Spectral confinement.} For every $i=2,\ldots,n$ and both signs,
$x_i^\sigma>-e/20$.
\item \emph{Local spectral mass.} For every vertex $p$, both signs, and every
$w\in[2e,2d]$,
\[
 \sum_{\substack{2\le i\le n\\x_i^\sigma\le w}}|v_i(p)|^2
 \le C_0w^{3/2}.
\]
\item \emph{Typical edge blocks.} For every pair
$w_B^\sigma\in[W_b,2W_b]\setminus\{x_i^\sigma:2\le i\le n\}$, define
\[
 R_B^\sigma=\sum_{\substack{2\le i\le n\\x_i^\sigma\ge w_B^\sigma}}
                 \frac{v_iv_i^T}{x_i^\sigma}.
\]
All but a fraction $C_0n^{-.1}$ of the $dn$ oriented edges $(p,p')$ satisfy
the following bound for both signs:
\[
 \left\|
 \begin{pmatrix}(R_B^\sigma)_{pp}&(R_B^\sigma)_{pp'}\\
                 (R_B^\sigma)_{p'p}&(R_B^\sigma)_{p'p'}\end{pmatrix}
 -\frac1{d-2}\begin{pmatrix}\sqrt q&\sigma\\\sigma&\sqrt q\end{pmatrix}
 \right\|_{\max}\le n^{-.024}.
\]
Here $\|\cdot\|_{\max}$ is the largest absolute entry.  The exceptional
edge set may depend on the chosen pair of cutoffs.
\end{enumerate}
Given the adjacency list of $G$, the algorithm finds $M$ vertex-disjoint
switches, each deleting two edges and inserting two edges.  In
$O_{d,C_0}(n^k)$ bit operations it outputs the adjacency list of a simple
$d$-regular graph $G'$ on the same vertices such that
\[
 -2\sqrt{d-1}<\lambda_i(G')<2\sqrt{d-1}
 \qquad (2\le i\le n).
\]
In particular, $G'$ is connected and nonbipartite.
\end{theorem}

\begin{proof}
Fix one admissible bulk cutoff for each sign.  We will bias the distribution
of an edge pair so that its expected switch improves the near spectrum at both
edges.  Second moments control the error from using only $M$ switches and from
eliminating the bulk.  Finally, we realize the draws in a polynomial finite
sample space and test every candidate exactly.

\par\smallskip\noindent\textit{Step 1: reduce positivity to the near subspace.}\quad
A draw consists of two oriented edges $g=(u,v)$ and $j=(x,y)$ of $G$.
With
$U_{g,j}=[e_u-e_x,e_v-e_y]$, replace $\{u,v\},\{x,y\}$ by
$\{u,y\},\{x,v\}$.  This is a valid switch when its four endpoints
are distinct and both inserted edges are absent from $G$.
Figure~\ref{fig:switch} illustrates the operation.  Degree-preserving
switches are classical in regular-graph sampling \cite{MW1990}.
Switching arguments for local Green-function comparison appear in
\cite[Sections~7 and~11]{BHY2019} and \cite[Sections~3 and~5]{HY2024};
the distribution here is chosen to improve both signed near-edge forms.
For a vertex-disjoint collection of $M$ valid switches,
concatenate their columns to obtain $U$.  The change of the signed precision is
\[
 sI-\sigma A_{G'}=sI-\sigma A_G+UC_\sigma U^T,
\]
where
\[
 C_\sigma=\operatorname{blockdiag}(\sigma F,\ldots,\sigma F),
 \qquad F=\begin{pmatrix}0&1\\1&0\end{pmatrix}.
\]
All columns of $U$ lie in $\mathbf 1^\perp$.
\begin{figure}[hbp]
\centering
\begin{tikzpicture}[scale=.9,
  vertex/.style={circle,fill=MidnightBlue,inner sep=2pt},
  edge/.style={thick,draw=MidnightBlue}]
\begin{scope}
\node[vertex,label=above:$u$] (u) at (0,1.5) {};
\node[vertex,label=above:$v$] (v) at (2,1.5) {};
\node[vertex,label=below:$x$] (x) at (0,0) {};
\node[vertex,label=below:$y$] (y) at (2,0) {};
\draw[edge] (u)--(v);
\draw[edge] (x)--(y);
\node at (1,-.8) {\small Original edges};
\end{scope}
\draw[-{Latex[length=2.5mm]},thick] (2.7,.75)--(4,.75);
\begin{scope}[xshift=4.7cm]
\node[vertex,label=above:$u$] (u2) at (0,1.5) {};
\node[vertex,label=above:$v$] (v2) at (2,1.5) {};
\node[vertex,label=below:$x$] (x2) at (0,0) {};
\node[vertex,label=below:$y$] (y2) at (2,0) {};
\draw[edge] (u2)--(y2);
\draw[edge] (x2)--(v2);
\node at (1,-.8) {\small Switched edges};
\end{scope}
\end{tikzpicture}
\caption{A degree-preserving switch deletes $uv$ and $xy$ and inserts
$uy$ and $xv$. The four endpoints are distinct, and the inserted
edges must be absent. The crossing in the right panel is not a vertex.
With $U=[e_u-e_x,e_v-e_y]$ and $F$ the two-coordinate swap matrix,
the adjacency change is $-UFU^T$.}
\label{fig:switch}
\end{figure}


For the fixed cutoffs, put
$E_0^\sigma=\operatorname{span}\{v_i:x_i^\sigma<w_B^\sigma\}$, and let
$U_E$ denote the coordinates of $U$ in this subspace.  Define
\begin{equation}
 S_B=C_\sigma^{-1}+U^TR_B^\sigma U.
 \label{eq:original-64}
\end{equation}
The two original edges in a draw each contribute one endpoint block to
$U^TR_B^\sigma U$.  Replacing these by their tree blocks and setting the
cross-edge entries to zero gives the ideal $2$-by-$2$ block
\[
 \sigma F+2P_{\mathrm{tree}}^\sigma
 =\frac1{d-2}
   \begin{pmatrix}s&\sigma d\\ \sigma d&s\end{pmatrix}.
\]
It has determinant $-1$, one eigenvalue of each sign, and inverse
\begin{equation}
 K_\sigma=
 \frac{d}{d-2}\begin{pmatrix}-s/d&\sigma\\ \sigma&-s/d\end{pmatrix}.
 \label{eq:original-65}
\end{equation}
If $S_B$ is sufficiently close in norm to the block diagonal matrix of these
ideal blocks, it has the same inertia as $C_\sigma^{-1}$.
To see why this suffices even when some old near slacks are negative, let
$D_B=\operatorname{diag}_{(E_0^\sigma)^\perp}(x_i^\sigma)>0$ be the
old bulk precision within $\mathbf1^\perp$, and let $U_B$ be the bulk
coordinates of $U$.  The two Schur complements of
\[
 \begin{pmatrix}D_B&U_B\\ U_B^T&-C_\sigma^{-1}\end{pmatrix}
\]
give
\[
 \operatorname{In}(D_B)+\operatorname{In}(-S_B)
 =\operatorname{In}(-C_\sigma^{-1})
   +\operatorname{In}(D_B+U_BC_\sigma U_B^T),
\]
where $\operatorname{In}$ records positive, negative and zero eigenvalue
counts.  The equal middle inertias cancel, proving that the updated bulk
precision $D_B+U_BC_\sigma U_B^T$ is positive definite.
Eliminating this bulk leaves the exact near Schur complement
\begin{equation}
 \operatorname{diag}_{E_0^\sigma}(x_i^\sigma)
       +U_ES_B^{-1}U_E^T.
 \label{eq:original-66}
\end{equation}
This is the Woodbury identity on the positive old bulk
\cite{HS1981}; no positivity of the old near block is required.

\par\smallskip\noindent\textit{Step 2: give the near spectrum a positive mean shift.}\quad
The inverse $K_\sigma$ in \eqref{eq:original-65} determines the ideal
contribution of one switch after the bulk has been eliminated.  Let
$U_{E,g,j}$ denote the near coordinates of the two columns of $U_{g,j}$,
and put
\[
 H^\sigma=U_{E,g,j}K_\sigma U_{E,g,j}^T.
\]
Thus replacing $S_B^{-1}$ in \eqref{eq:original-66} by the block diagonal
matrix with blocks $K_\sigma$ replaces the near correction by
$\sum_{r=1}^M H_r^\sigma$.  Since $K_\sigma$ is indefinite, we must control
this sum through the choice of the edge-pair distribution.

First let $\nu_0$ be the law of two independent uniformly chosen oriented
edges.  For $g=(u,v)$ and nonconstant indices $i$, write
\[
 z_i(g)=(v_i(u),v_i(v))^T,\qquad
 M_i=\begin{pmatrix}1&\lambda_i/d\\ \lambda_i/d&1\end{pmatrix}.
\]
Regularity and orthogonality give
\[
 \mathbb E_{\nu_0}z_i(g)=0,\qquad
 \mathbb E_{\nu_0}[z_i(g)z_k(g)^T]=\frac{\delta_{ik}}n M_i.
\]
For indices in $E_0^\sigma$, the entry of the single-switch correction is
\[
 H_{ik}^\sigma=(z_i(g)-z_i(j))^TK_\sigma(z_k(g)-z_k(j)).
\]
Its uniform mean is
\[
 \mathbb E_{\nu_0}H_{ik}^\sigma
 =\frac{2\delta_{ik}}n\operatorname{Tr}(K_\sigma M_i)
 =-\frac{4\delta_{ik}x_i^\sigma}{(d-2)n}.
\]
Consequently the uniform mean correction only multiplies each old slack by
$1-4M/((d-2)n)>0$.  It cannot turn a negative slack into a positive one.
We therefore bias the correlations between the two edges on the modes
closest to either threshold.

Choose cutoffs $t_T^\tau\in[2e,3e]$ avoiding the signed slacks
$\{x_i^\tau:2\le i\le n\}$, and set
\[
 T_\tau=\{i:x_i^\tau<t_T^\tau\},\qquad
 J_T\asymp n^{3a/2+a/10}.
\]
Define the tilted law $\nu$ by multiplying the $\nu_0$ probabilities by
\begin{equation}
 1+\frac1{B_*}\left[-\frac{2n}{J_T}
       \sum_{\tau=\pm1}\sum_{i\in T_\tau}
       z_i(g)^TK_\tau z_i(j)\right],
 \label{eq:original-67}
\end{equation}
where $B_*>0$ is a sufficiently large fixed constant.
This is a positive probability density.  Indeed,
\eqref{eq:original-61} and Cauchy--Schwarz bound the absolute value of the
bracket by
\[
 C\frac{n e^{3/2}}{J_T}=O(n^{-a/10})=o(1).
\]
Every nonconstant eigenvector has mean zero on either endpoint of a uniform
oriented edge.  The bracket therefore has mean zero under $\nu_0$ even when
either edge is fixed.  The density is normalized and both edge marginals
under $\nu$ remain uniform.

Expansion under this tilted law gives
\begin{equation}
 \mathbb E_\nu H_{ik}^\sigma
 =\delta_{ik}\left[
 -\frac{4x_i^\sigma}{(d-2)n}
 +\frac4{B_*J_Tn}
   \sum_{\tau:\,i\in T_\tau}
      \operatorname{Tr}(K_\sigma M_iK_\tau M_i)
 \right].
 \label{eq:original-68}
\end{equation}
The first term is the uniform mean computed above.  In the tilt correction
only the two cross-edge terms in the expansion of $H_{ik}^\sigma$ survive.
Each forces the two near indices to agree with the tilting index, and each
contributes the same trace divided by $n^2$.  Their negative signs cancel
the negative sign in \eqref{eq:original-67}, giving the displayed factor
$4/(B_*J_Tn)$.

For a $\sigma$-near mode, only $T_\sigma$ can occur in the sum, since the two
signed slacks add to $2s$ and $2W_b+3e=o(1)$.
On these modes $M_i$ is uniformly positive definite, because $s<d$, and
\[
 \operatorname{Tr}(K_\sigma M_iK_\sigma M_i)
 =\|M_i^{1/2}K_\sigma M_i^{1/2}\|_{\mathrm F}^2\ge c_d.
\]
The positive bias satisfies
$M/(J_Tn)\asymp n^{-2/3+8.4a}\gg e$.
It overwhelms the possible slack $-e/20$ on $T_\sigma$.
Outside $T_\sigma$, the slack is at least $2e$, and multiplication by
$1-4M/((d-2)n)=1-o(1)$ preserves a fixed fraction of it.  Thus the same
law improves both signed forms:
\begin{equation}
 \operatorname{diag}(x_i^\sigma)+M\mathbb E_\nu H^\sigma\succeq ceI
 \quad\text{on }E_0^\sigma.
 \label{eq:original-69}
\end{equation}

\par\smallskip\noindent\textit{Step 3: restrict the draws to typical edges.}\quad
For the bulk cutoffs fixed in Step 1, keep only oriented edges whose endpoint
blocks have the required accuracy for both signs.
The numerical implementation below includes every edge with the proven
$n^{-.024}$ accuracy and includes only edges with $Cn^{-.02}$ accuracy.
Condition $\nu$ on both sampled edges lying in this set.  The removed
probability is $O(n^{-.1})$.  Approximate the resulting law in total variation
by $n^{-C}$, for a sufficiently large fixed $C$, retaining a bound
$C_d/(dn)^2$ for every point probability.

For any draw, the projector estimate at scale $2W_b$ gives
\[
 \|U_{E,g,j}\|_{\mathrm F}^2\le CW_b^{3/2},\qquad
 \|H^\sigma\|_{\mathrm F}\le CW_b^{3/2}.
\]
Consequently, the restriction and the probability approximation change the
left side of \eqref{eq:original-69} by at most
\[
 CM(n^{-.1}+n^{-C})W_b^{3/2}=o(e).
\]
Denote the approximating law, including the probability rounding in Step~7,
by $\mu$.  Unlike $\nu$, it need not have exactly uniform marginals; the
pointwise density bound is what the remaining estimates use.

\par\smallskip\noindent\textit{Step 4: control fluctuations and collisions.}\quad
Take $M$ draws with law $\mu$, requiring only pairwise independence between
draws.  The bulk inverse error will be sandwiched by $U_E$, so a bound
$\|U_EU_E^T\|=O(M/n)$ makes $o(en/M)$ the bulk error scale needed to
obtain an $o(e)$ error in the near form.  We prove that, with probability
$1-o(1)$, simultaneously for both signs,
\begin{equation}
 \begin{split}
 &\left\|\sum_{r=1}^M H_r^\sigma-M\mathbb E_\mu H^\sigma\right\|=o(e),\\
 &\|U_EU_E^T\|\le CM/n,\\
 &\left\|S_B-\operatorname{blockdiag}
          (\sigma F+2P_{\mathrm{tree}}^\sigma)\right\|=o(en/M).
 \end{split}
 \label{eq:original-70}
\end{equation}
Here and below a block diagonal notation in this display has $M$ blocks.
We supply all three second-moment estimates.

Pairwise independence makes the cross terms in the squared Frobenius norm of
a centered sum vanish.  Hence
\[
 \mathbb E\left\|\sum_{r=1}^M
 (H_r^\sigma-\mathbb E_\mu H_r^\sigma)\right\|_{\mathrm F}^2
 \le CM W_b^3.
\]
For the sum $U_EU_E^T$ the same estimate applies to its centered version.
Uniform-density domination gives its mean the positive-order bound $CMI/n$:
for any unit vector in the near subspace, average the squares of its four
endpoint values and use that an endpoint of a uniform oriented edge is uniform
on vertices.  Since
\[
 \sqrt M\,W_b^{3/2}=n^{-5/6+155a+o(1)}
       \ll e,\ M/n,
\]
Markov's inequality proves the first two assertions.

To prove the third, first retain the two endpoint blocks belonging to each
separate sampled edge.  They have error $O(n^{-.02})$ by construction; after
forming the difference columns this is still a block-diagonal operator error
$O(n^{-.02})$.
All the other entries concern endpoints in different edge slots.  This
includes the two edge slots within a draw.  Such a pair of edge slots has joint
density bounded by a constant relative to two uniform oriented edges: within
one draw this is the pointwise bound for $\mu$, and between draws it follows
from pairwise independence and the bounded marginal densities.
For any prescribed endpoint in each slot, \eqref{eq:original-63} therefore
gives
\[
 \mathbb E\left| (R_B^\sigma)_{pq}\right|^2
 \le \frac{C}{n^2}\sum_{p,q}|(R_B^\sigma)_{pq}|^2
 \le \frac{C}{n\sqrt{W_b}}.
\]
There are $O(M^2)$ such entries.  Passing from the endpoint columns to their
differences has operator norm $\sqrt2$, so the total squared Frobenius budget
is at most $CM^2/(n\sqrt{W_b})$ up to a fixed factor.  Its square root obeys
\[
 \frac{M}{\sqrt n\,W_b^{1/4}}
   =n^{b_*-t_*/4+o(1)}=n^{-15a+o(1)}
   \ll en/M=n^{-9a+o(1)}.
\]
This is exactly the strict inequality $t_*>8b_*-4a$.
Also $n^{-.02}=o(en/M)$.  Markov's inequality now proves the third assertion
of \eqref{eq:original-70}.  All the displayed separations have a fixed power
margin, so the $o(1)$ probability bounds can be imposed simultaneously.

The same density estimates show that endpoint coincidences within or between
draws, or a proposed new edge already present in $G$, have total probability
$O_d(M^2/n)=n^{-1/3+20a+o(1)}=o(1)$.
Indeed each particular coincidence or adjacency has probability $O_d(1/n)$;
there are $O(M^2)$ pairs of slots.  Original edges have distinct endpoints by
simplicity.  On the complementary event all four endpoints of every switch
are distinct and different switches are vertex-disjoint.  Thus the switches
are simultaneously valid, introduce neither loops nor repeated edges, and
preserve every degree.

The following ratios collect the margins used above.  Fixed factors and
integer rounding do not change their powers of $n$.  The first ratio grows;
every error ratio tends to zero for $a=10^{-4}$.
\begin{center}
\small
\begin{tabular}{@{}lll@{}}
\toprule
Contribution&Comparison scale&Ratio, up to fixed factors\\
\midrule
Mean bias $M/(J_Tn)$&$e$&$n^{7.4a}$\\
Conditioning error $Mn^{-.1}W_b^{3/2}$&$e$&$n^{-.1+159a}$\\
Near fluctuation $\sqrt M\,W_b^{3/2}$&$e$&$n^{-1/6+154a}$\\
Bulk fluctuation $M/(\sqrt n\,W_b^{1/4})$&$en/M$&$n^{-6a}$\\
Typical-block error $n^{-.02}$&$en/M$&$n^{-.02+9a}$\\
Collision probability $M^2/n$&$1$&$n^{-1/3+20a}$\\
\bottomrule
\end{tabular}
\end{center}
Since $e=o(M/n)$, the near-fluctuation estimate also supplies the
second line of \eqref{eq:original-70}.  No common exceptional edge set over
different cutoff pairs is used: the cutoffs were fixed before the law
was conditioned.

\par\smallskip\noindent\textit{Step 5: conclude strict positivity at both edges.}\quad
For a collection satisfying all these conclusions, the inverse perturbation
identity and the bounded inverse in \eqref{eq:original-65} give
\[
 \left\|S_B^{-1}-\operatorname{blockdiag}(K_\sigma)\right\|
 =o(en/M).
\]
Write $\mathcal K_\sigma=\operatorname{blockdiag}(K_\sigma,\ldots,K_\sigma)$.
The near Schur complement is exactly
\[
 \operatorname{diag}(x_i^\sigma)+\sum_{r=1}^M H_r^\sigma
       +U_E(S_B^{-1}-\mathcal K_\sigma)U_E^T.
\]
The last term has norm $o(e)$ by the second line of
\eqref{eq:original-70}.  The centered fluctuation and the change from
$\nu$ to $\mu$ each cost another $o(e)$.  The positive mean bound
\eqref{eq:original-69} therefore makes this Schur complement positive
definite.
The updated bulk was already positive definite.  Thus the full signed
precision is positive definite on $\mathbf 1^\perp$, for both signs.
This proves the asserted strict spectral bounds for at least one collection.

\par\smallskip\noindent\textit{Step 6: compute rational approximations to the law.}\quad
It remains to find such a collection deterministically.  All accuracies in
this paragraph can be $n^{-C}$ for any sufficiently large fixed $C$.
Use the rational matrix
\[
 A_\perp=A_G-d\mathbf 1\mathbf 1^T/n,
\]
which agrees with $A_G$ on $\mathbf 1^\perp$ and has eigenvalue zero on the
constant direction.  For each sign, choose rational cutoffs for the eigenvalues
of $\sigma A_\perp$, with the corresponding slack thresholds in the interiors
of $[W_b,2W_b]$ and $[2e,3e]$, respectively.  These cutoffs can have a gap at
least an inverse polynomial from the spectrum.  To see this, approximate $s$
to much greater precision, place more than $2n$ rational candidates at spacing
comparable to the interval width divided by $n$, and put disjoint small
neighborhoods around them.  At most $n$ neighborhoods can contain eigenvalues.
Exact rational inertia tests at the endpoints, together with a zero-eigenvalue
test at an endpoint, identify an empty neighborhood.  Choosing the candidates
in a fixed interior subinterval ensures that the error in approximating $s$
does not change the required slack interval.

The projectors onto $T_\tau$ and the operators $R_B^\sigma$ can now be
approximated in operator norm without resolving any internal eigenvalue
spacings.  On a fixed rational interval containing $[-d,d]$, define a continuous
cutoff function equal to its required zero or one value outside the separating
gap, with linear interpolation inside the gap.  For the bulk reciprocal use a
very accurate rational $\widetilde s$: set the function equal to
$1/(\widetilde s-t)$ on the bulk side and zero on the excluded side, and
interpolate through the empty gap.  These functions have polynomially bounded
Lipschitz constants and rational values of polynomial bit length at rational
arguments.  The reciprocal approximation has an additional error at most
$C|\widetilde s-s|/W_b^2$, which is made negligible by the chosen precision.

After affine scaling to $[0,1]$, a function with Lipschitz constant $L_f$ has
Bernstein polynomial \cite{Bernstein1912} of degree $D$ with uniform error at most
$L_f/(2\sqrt D)$.  This follows by writing the Bernstein polynomial as
$\mathbb E f(X/D)$ for $X\sim\operatorname{Bin}(D,t)$ and applying
Cauchy--Schwarz to $\mathbb E|X/D-t|$.
Take a polynomially large degree $D$ to obtain the desired inverse-polynomial
accuracy, and evaluate these rational polynomials at $\sigma A_\perp$.
The spectral theorem transfers the scalar error to operator norm.
Sandwich the resulting operators by $I-\mathbf 1\mathbf 1^T/n$ to remove the
artificial constant eigenvalue.  Rational matrix-polynomial evaluation has
polynomial bit complexity: the degree is polynomial, its coefficients have
polynomial bit length, and direct products and sums increase bit lengths by
at most polynomial factors.

Use the approximate bulk blocks and approximate
$P_{\mathrm{tree}}^\sigma$ to accept all edges with the proved
$n^{-.024}$ accuracy and only edges with $Cn^{-.02}$ accuracy.  For example,
use a rational entrywise threshold between these two scales; the approximation
error is much smaller than either threshold.  The sum in the tilt is
basis-invariant: for endpoint indices $g_r,j_t$ it equals
\[
 \sum_{r,t=1}^2(K_\tau)_{rt}
       (\operatorname{Proj}_{T_\tau})_{g_r,j_t}.
\]
Thus the approximate projectors and rational approximations to $K_\tau$
give rational positive weights for every accepted edge pair.  Normalize these
weights.  Errors can be made small enough, including the factor $n/J_T$, that
the law is within $n^{-C}$ in total variation of the conditional tilted law.
Its point probabilities remain bounded by $C_d/(dn)^2$.

\par\smallskip\noindent\textit{Step 7: enumerate a polynomial sample space.}\quad
Choose a prime $P$ between $n^{C'}$ and $2n^{C'}$, up to immaterial rounding,
where $C'$ is a sufficiently large fixed constant and $P>M$.
Bertrand's theorem and deterministic trial search provide such a prime in
time polynomial in $n$.
There are $O_d(n^2)$ accepted ordered edge-pair outcomes.  Floor their
probabilities times $P$, then distribute the remaining counts among the
largest fractional remainders.  The counts are nonnegative, sum to $P$,
and differ by at most one from their unrounded values; outcomes outside
the accepted support retain count zero.  Assign that many residues
of $\mathbb F_P$ to each outcome.  The total variation error is
$O_d(n^2/P)$ and the maximum probability increases by at most $1/P$; both are
within the stipulated bounds for large enough $C'$.
Choose the preceding approximation precision and $C'$ so that their combined
total-variation error is at most $n^{-C}$.  The resulting rounded law is a
choice of $\mu$ with all the properties used in Steps 3--5.

Use the standard affine construction of pairwise-independent variables
(see \cite[Section~2.2]{PR2008}).  Enumerate all $P^2$ seeds
$(\gamma,\beta)\in\mathbb F_P^2$ and use the residues
\[
 \gamma r+\beta\qquad(1\le r\le M)
\]
to select the draws.  Under a uniform seed, any two residues are independent
and uniform: for $r\ne t$, the map
$(\gamma,\beta)\mapsto(\gamma r+\beta,\gamma t+\beta)$ is bijective since
$r-t\ne0$ in $\mathbb F_P$.  Mapping residues to outcomes preserves pairwise
independence.  The estimates above therefore prove that a positive fraction,
indeed a fraction $1-o(1)$, of these seeds satisfy all the required events.

For each seed, discard collections with intersecting endpoint sets or an
invalid switch, build the candidate graph, and test its spectrum exactly.
A rational contrast basis tests the positivity
of
\[
 4qI-A_{G'}^2\quad\text{on }\mathbf 1^\perp
\]
by exact rational congruence and inertia.  Positive definiteness gives the
strict conclusion of the theorem.  Symmetric elimination uses a nonzero
diagonal pivot if available, otherwise an invertible off-diagonal
$2$-by-$2$ pivot, with a remaining zero block contributing zero eigenvalues.
Fraction-free arithmetic or determinant bounds keep all intermediate bit
lengths polynomial: Schur-complement entries are ratios of minors of the
original pivoted rational matrix.  The same procedure implements the cutoff
tests above.  A passing seed exists by the preceding argument, and the number
of seeds and all work per seed are polynomial in $n$.

Finally, $d>2\sqrt{d-1}$.  A disconnected regular graph has an additional
$d$-eigenvector orthogonal to $\mathbf 1$, and a bipartite regular graph has a
$-d$-eigenvector.  The verified strict bound excludes both possibilities.
\end{proof}

\begin{remark}[Inputs supplied by the pairing construction]
The theorem uses only the three stated properties of $G$.
The completed pairing supplies $x_i^\sigma>-e/20$ and the final entry
comparison.  Its frozen-reference conclusions verify
Lemma~\ref{lem:repair-projector}; its separation-stop path estimates verify
the typical-edge assertion in Lemma~\ref{lem:repair-bulk}.  The repair therefore
applies to precisely the graph constructed in the preceding sections, with
no change of vertex set.
\end{remark}
